%% Fig 1: Conceptual landscape mindmap (8 research domains)
%% Uses TikZ mindmap library with Okabe-Ito colors
\begin{tikzpicture}[
  mindmap,
  every node/.style={concept, circular drop shadow, minimum size=0pt},
  root concept/.append style={
    concept color=black!40, fill=white, line width=1.2pt,
    font=\sffamily\bfseries\large, text=black,
    minimum size=3.5cm,
  },
  level 1 concept/.append style={
    font=\sffamily\small\bfseries, text=white,
    minimum size=2.2cm, level distance=4.5cm,
    sibling angle=45,
  },
  level 2 concept/.append style={
    font=\sffamily\tiny, text=white,
    minimum size=1.2cm, level distance=2.8cm,
  },
]
\node[root concept] {Algebraic\\Structure\\Constrains\\Dynamics}
  child[concept color=cAlgebra, grow=90] {
    node[concept] {Algebra}
    child[grow=60] { node[concept] {Cayley--\\Dickson} }
    child[grow=120] { node[concept] {Box-\\Kites} }
  }
  child[concept color=cFluid, grow=45] {
    node[concept] {Fluid\\Dynamics}
    child[grow=30] { node[concept] {LBM\\D3Q19} }
  }
  child[concept color=cGR, grow=0] {
    node[concept] {General\\Relativity}
    child[grow=-20] { node[concept] {Warp\\Bubble} }
  }
  child[concept color=cQuantum, grow=-45] {
    node[concept] {Quantum}
    child[grow=-60] { node[concept] {CHSH\\Bridge} }
  }
  child[concept color=cMaterials, grow=-90] {
    node[concept] {Materials}
    child[grow=-110] { node[concept] {Meta-\\material} }
  }
  child[concept color=cOptics, grow=-135] {
    node[concept] {Optics}
    child[grow=-150] { node[concept] {Casimir} }
  }
  child[concept color=cStatistics, grow=180] {
    node[concept] {Statistics}
    child[grow=200] { node[concept] {Ghost\\Audit} }
  }
  child[concept color=cCosmology, grow=135] {
    node[concept] {Cosmology}
    child[grow=150] { node[concept] {Dark\\Energy} }
  };
\end{tikzpicture}
